% © 2026 Ing. David Jaroš — CC BY-NC-ND 4.0
%
% This work is licensed under a Creative Commons Attribution-NonCommercial-NoDerivatives
% 4.0 International License (CC BY-NC-ND 4.0).
%
% License History: Earlier drafts (up to v0.3) were released under CC BY 4.0.
% From v0.4 onward, all material is released under CC BY-NC-ND 4.0 to protect
% the integrity of the theoretical work during ongoing academic development.

\documentclass[11pt,a4paper]{article}
\usepackage{amsmath, amssymb, amsthm}
\usepackage{hyperref}
\usepackage{geometry}
\geometry{margin=2.5cm}

\newtheorem{theorem}{Theorem}
\newtheorem{proposition}{Proposition}
\newtheorem{definition}{Definition}
\newtheorem{remark}{Remark}
\newtheorem{conjecture}{Conjecture}

\title{Modular Dynamics in Unified Biquaternion Theory\\
\large Torus Construction, Modular Invariants, and Fixing \texorpdfstring{$\mathrm{Im}(\tau)$}{Im(tau)}}
\author{Ing.\ David Jaroš}
\date{2026}

\begin{document}
\maketitle

\begin{abstract}
We continue the analysis of modular structures connecting the complex time
$\tau$, the torus geometry $\mathbb{C}/(\mathbb{Z}+\tau\mathbb{Z})$, theta
functions, and Hecke operators.  Building on the torus construction of
\texttt{research/theta\_modular\_geometry.tex} and the moduli-space analysis
of \texttt{research/moduli\_space\_ads\_vs\_physical\_ds.tex}, we examine
candidate mechanisms for fixing $\mathrm{Im}(\tau) = R_\psi/R_t$ from the
UBT field dynamics.  All claims are labelled with their proof status.
\end{abstract}

\tableofcontents

% ======================================================================
\section{Setup: The Modular Parameter \texorpdfstring{$\tau$}{tau} in UBT}
\label{sec:tau}
% ======================================================================

\subsection{Recap of the torus construction}

Following \texttt{research/theta\_modular\_geometry.tex} §1--2, the complex
time $\tau = t + i\psi$ with $\psi \sim \psi + 2\pi R_\psi$ defines a
modular parameter:
\begin{equation}
  \tau_{\mathrm{mod}} = i\frac{R_\psi}{R_t} \in \mathbb{H},
\end{equation}
lying on the imaginary axis of the upper half-plane $\mathbb{H}$.  The
associated torus is:
\begin{equation}
  \mathbb{T}_\tau = \mathbb{C}/(\mathbb{Z} + \tau_{\mathrm{mod}}\mathbb{Z}).
\end{equation}

\begin{remark}[Status of the construction]
  The torus construction is a proved consequence of the axioms (Status:
  \textbf{[L1]}).  The identification $\tau_{\mathrm{mod}} = iR_\psi/R_t$
  places $\tau_{\mathrm{mod}}$ on the imaginary axis, which is a special
  locus in the moduli space.  See
  \texttt{AUDITS/complex\_time\_modular\_audit.md}.
\end{remark}

% ======================================================================
\section{Modular Invariants}
\label{sec:modular_inv}
% ======================================================================

\subsection{Physical quantities and modular invariance}

Observable quantities (particle masses, coupling constants) must be
modular invariants (or covariant under the relevant subgroup of
$\mathrm{SL}(2,\mathbb{Z})$).

The modular $j$-invariant:
\begin{equation}
\label{eq:j_inv}
  j(\tau) = 1728\, \frac{g_2(\tau)^3}{\Delta(\tau)},
\end{equation}
where $g_2, \Delta$ are the Eisenstein series and discriminant, is the
prototypical invariant.

\begin{remark}[Status of modular invariance in UBT]
  The UBT action $S[\Theta]$ has not been shown to be fully modular invariant.
  The kinetic term is invariant under the continuous scaling $R_\psi \to \lambda R_\psi$
  (Weyl symmetry), but invariance under the full $\mathrm{SL}(2,\mathbb{Z})$
  transformation $\tau \to (a\tau+b)/(c\tau+d)$ has not been checked.

  Status: \textbf{[O]} — open problem.  Establishing modular invariance
  would be a major structural result for UBT.
\end{remark}

% ======================================================================
\section{Candidate Mechanisms for Fixing \texorpdfstring{$\mathrm{Im}(\tau)$}{Im(tau)}}
\label{sec:fix_tau}
% ======================================================================

\subsection{Mechanism M1: Moduli potential from the Θ-field action}

A moduli potential $V_{\mathrm{mod}}(\tau)$ can arise from the one-loop
determinant of the Θ-field:
\begin{equation}
\label{eq:moduli_potential}
  V_{\mathrm{mod}}(\tau)
  = -\ln\det(-\Delta_{T^3}(\tau))
  = -\ln\prod_{n\in\mathbb{Z}^3\setminus\{0\}}
    \frac{|\mathbf{n}|^2}{R_\psi^2(\tau)}.
\end{equation}

This potential extremises at special values of $\tau$.  For a product torus,
the minimum of $V_{\mathrm{mod}}$ is at the self-dual point $R_\psi = R_t$,
i.e., $\tau_{\mathrm{mod}} = i$.

\begin{remark}[Status]
  The existence of a non-trivial moduli potential from the one-loop $\Theta$
  determinant is a standard result in string theory and Kaluza--Klein
  compactifications.  Applied to UBT, it predicts $R_\psi \sim R_t$
  (no hierarchy between real and imaginary time scales) unless additional
  contributions shift the minimum.  Status: \textbf{[L2]} — motivated by
  analogy, not computed for UBT.
\end{remark}

\subsection{Mechanism M2: Eisenstein series and enhanced symmetry points}

The modular $j$-function has special values at:
\begin{equation}
  j(i) = 1728, \quad j(e^{2\pi i/3}) = 0.
\end{equation}

These correspond to tori with automorphism groups $\mathbb{Z}_4$ and
$\mathbb{Z}_6$ respectively.  If the UBT partition function is invariant
under the automorphism group of the torus, then $\tau_{\mathrm{mod}}$
may be stabilised at these special points.

\begin{conjecture}[Stabilisation at $j = 1728$]\label{conj:j1728}
  The modular potential $V_{\mathrm{mod}}$ has a minimum at $\tau_{\mathrm{mod}} = i$,
  corresponding to $R_\psi = R_t$ and $j = 1728$.  This fixes
  $\mathrm{Im}(\tau_{\mathrm{mod}}) = 1$.

  \textbf{Status: [L2]}.  The conjecture is motivated by the enhanced
  $\mathbb{Z}_4$ automorphism at $\tau = i$, but no computation of
  $V_{\mathrm{mod}}$ from the UBT action has been performed.
\end{conjecture}

\subsection{Mechanism M3: Theta function zeros and spectral gaps}

The torus theta function $\vartheta_3(0|\tau_{\mathrm{mod}})$ has a zero
at $\tau_{\mathrm{mod}} = (1 + i)/2$ (boundary of the fundamental domain).
A spectral gap opens when $\vartheta_3$ passes through zero, which could
stabilise $\tau_{\mathrm{mod}}$ near this value.

\begin{remark}[Status]
  This mechanism is speculative.  Status: \textbf{[L2]}.
\end{remark}

% ======================================================================
\section{Modular Forms and Their Connection to UBT}
\label{sec:mf}
% ======================================================================

\subsection{Theta series from KK spectrum}

The partition function of the KK tower on $T^3$ is:
\begin{equation}
\label{eq:kk_partition}
  Z_{\mathrm{KK}}(\tau)
  = \sum_{\mathbf{n}\in\mathbb{Z}^3}
    q^{|\mathbf{n}|^2},
  \quad q = e^{2\pi i\tau_{\mathrm{mod}}},
\end{equation}
which is the level-1 theta series $\vartheta_3^3(\tau)$.

\begin{remark}[Modular transformation]
  Under $\tau \to \tau + 1$: $q \to e^{2\pi i}q = q$, so
  $Z_{\mathrm{KK}}$ is invariant.
  
  Under $\tau \to -1/\tau$: Poisson duality (see
  \texttt{research/B\_base\_spectral\_determinant.tex} §4) gives
  $Z_{\mathrm{KK}}(-1/\tau) = |\tau|^{3/2} Z_{\mathrm{KK}}(\tau)$.

  The $\tau \to -1/\tau$ transformation is \emph{not} a symmetry of
  $Z_{\mathrm{KK}}$; it changes it by the factor $|\tau|^{3/2}$.
  Full modular invariance would require a specific combination of
  winding and momentum modes, as in string theory.
  Status: \textbf{[L0]}.
\end{remark}

% ======================================================================
\section{Connection to the \texorpdfstring{$B_{\mathrm{base}}$}{B\_base} Programme}
\label{sec:bbase}
% ======================================================================

If the modular potential fixes $\tau_{\mathrm{mod}} = i$ (Conjecture~\ref{conj:j1728}),
then $R_\psi = R_t$ and the $N_{\mathrm{eff}}$ parameter becomes:
\begin{equation}
  N_{\mathrm{eff}} = [\det(-\Delta_{T^3}(i))]^{1/3}.
\end{equation}

At $\tau_{\mathrm{mod}} = i$, the Epstein zeta function takes the value:
\begin{equation}
  Z_3(s)|_{\tau=i} = 6\,\zeta(2s)\,\beta(2s-1),
\end{equation}
where $\beta$ is the Dirichlet $\beta$-function.  This gives a definite
numerical value for $N_{\mathrm{eff}}$, potentially connecting to
$B_{\mathrm{base}} = 41.57$.

\begin{remark}[Status of the connection]
  Computing $N_{\mathrm{eff}}$ from $Z_3(s)|_{\tau=i}$ and checking
  whether $N_{\mathrm{eff}}^{3/2} = 41.57$ is a concrete numerical task
  that has not yet been performed.  Status: \textbf{[O]}.
  This is a high-priority next step in the $B_{\mathrm{base}}$ programme.
\end{remark}

% ======================================================================
\section{Open Problems}
\label{sec:open}
% ======================================================================

\begin{enumerate}
  \item \textbf{[O] Compute $V_{\mathrm{mod}}(\tau)$}: Evaluate the one-loop
    moduli potential from the Θ-field determinant and find its extrema.

  \item \textbf{[O] Modular invariance of the UBT action}: Check whether
    $S[\Theta]$ is invariant under $\tau \to (a\tau+b)/(c\tau+d)$.

  \item \textbf{[O] $N_{\mathrm{eff}}$ at $\tau = i$}: Compute
    $N_{\mathrm{eff}} = [\det(-\Delta_{T^3}(i))]^{1/3}$ numerically.

  \item \textbf{[O] Hecke operators and generation structure}: Connect
    the Hecke eigenvalues discussed in
    \texttt{research/hecke\_generation\_structure.tex} to the modular
    structure of the Θ-field partition function.
\end{enumerate}

% ======================================================================
\section{Cross-References}
\label{sec:xref}
% ======================================================================

\begin{tabular}{ll}
\hline
Topic & File \\
\hline
Torus construction & \texttt{research/theta\_modular\_geometry.tex} \\
Moduli space geometry & \texttt{research/moduli\_space\_ads\_vs\_physical\_ds.tex} \\
Modular audit & \texttt{AUDITS/complex\_time\_modular\_audit.md} \\
Poisson duality & \texttt{research/B\_base\_spectral\_determinant.tex} §4 \\
Heat kernel & \texttt{research/B\_base\_heat\_kernel.tex} \\
Hecke structure & \texttt{research/hecke\_generation\_structure.tex} \\
\hline
\end{tabular}

\end{document}
